% GENERATED FILE. Edit canonical lessons, then run npm run generate:books.
% canonical-source-hash: fnv1a64:bbd8a7fea8e281e3
% canonical-lessons: RU-C491-svekrov, RU-C491-tyoshcha, RU-C491-svyokor, RU-C491-test, RU-C491-zyat

\chapter{In-laws}
\label{ch:ru-in-laws}
\hlchaptermodality{491}

\begin{chapteropening}
\textbf{By the end of this chapter:} \emph{I can talk about my in-laws.}

Complete the last lesson of chapter 491: I can talk about my in-laws.
\end{chapteropening}

\section[\texorpdfstring{\ru{свекровь} (svekróv)}{свекровь (svekróv)}]{\ru{свекровь} (svekróv) — a husband's mother, a mother-in-law}

\label{lesson:RU-C491-svekrov}



\begin{quote}
Before the new one: say the Russian for to swallow, then the Russian for syrup.
\end{quote}

\subsection*{You'll want to know: \ru{свекровь}}
\textbf{\ru{свекровь}} (\emph{svekróv}) — \textquotedblleft{}a husband's mother, a mother-in-law\textquotedblright{}.

\begin{quote}
\ru{Моя} \ru{свекровь} \ru{прекрасно} \ru{готовит}.
\end{quote}

\begin{quote}
\emph{Moyá svekróv prekrásno gotóvit.}
\end{quote}

\begin{quote}
My mother-in-law is a wonderful cook.
\end{quote}

Feminine. The husband's mother.

\begin{cousinweb}
\emph{\ru{свекровь}} is inherited from Proto-Slavic \emph{svekry}, which EtymDB traces to PIE \emph{swe\'{k}rúh\textsubscript{2} 'mother-in-law'.}
\end{cousinweb}

\begin{grammarlens}[title={What you've built}]
One more word for telling, explaining and describing at B1.
\end{grammarlens}

\subsection*{Your turn}
\begin{itemize}
\raggedright
  \item \emph{Say it:} \emph{svekróv}
  \item \emph{Say it:} \emph{svekróv}, once more
  \item \emph{Say it:} a sentence of your own with \emph{svekróv}
  \item \emph{Recall:} say the Russian for to swallow, then the Russian for syrup, then say \emph{svekróv} again
\end{itemize}

\begin{tcolorbox}[breakable,colback=teal!4,colframe=teal!35!black,title={Before you move on}]
What does \ru{свекровь} mean? (\textquotedblleft{}A husband's mother, a mother-in-law\textquotedblright{}.) Say it once more.
\end{tcolorbox}
\section[\texorpdfstring{\ru{тёща} (tyóshcha)}{тёща (tyóshcha)}]{\ru{тёща} (tyóshcha) — a wife's mother, a mother-in-law}

\label{lesson:RU-C491-tyoshcha}



\begin{quote}
Before the new one: say the Russian for syrup, then the Russian for a husband's mother.
\end{quote}

\subsection*{You'll want to know: \ru{тёща}}
\textbf{\ru{тёща}} (\emph{tyóshcha}) — \textquotedblleft{}a wife's mother, a mother-in-law\textquotedblright{}.

\begin{quote}
\ru{Тёща} \ru{часто} \ru{приезжает} \ru{к} \ru{нам} \ru{на} \ru{выходные}.
\end{quote}

\begin{quote}
\emph{Tyóshcha chásto priyezzháyet k nam na vykhodnýye.}
\end{quote}

\begin{quote}
My mother-in-law often comes to us for the weekend.
\end{quote}

Feminine. The wife's mother.

\begin{cousinweb}
EtymDB has no entry for \emph{\ru{тёща}}; it is often linked to \emph{\ru{тесть}} 'wife's father', but the origin is not settled.
\end{cousinweb}

\begin{grammarlens}[title={What you've built}]
One more word for telling, explaining and describing at B1.
\end{grammarlens}

\subsection*{Your turn}
\begin{itemize}
\raggedright
  \item \emph{Say it:} \emph{tyóshcha}
  \item \emph{Say it:} \emph{tyóshcha}, once more
  \item \emph{Say it:} a sentence of your own with \emph{tyóshcha}
  \item \emph{Recall:} say the Russian for syrup, then the Russian for a husband's mother, then say \emph{tyóshcha} again
\end{itemize}

\begin{tcolorbox}[breakable,colback=teal!4,colframe=teal!35!black,title={Before you move on}]
What does \ru{тёща} mean? (\textquotedblleft{}A wife's mother, a mother-in-law\textquotedblright{}.) Say it once more.
\end{tcolorbox}
\section[\texorpdfstring{\ru{свёкор} (svyókor)}{свёкор (svyókor)}]{\ru{свёкор} (svyókor) — a husband's father, a father-in-law}

\label{lesson:RU-C491-svyokor}



\begin{quote}
Before the new one: say the Russian for a husband's mother, then the Russian for a wife's mother.
\end{quote}

\subsection*{You'll want to know: \ru{свёкор}}
\textbf{\ru{свёкор}} (\emph{svyókor}) — \textquotedblleft{}a husband's father, a father-in-law\textquotedblright{}.

\begin{quote}
\ru{Свёкор} \ru{помог} \ru{нам} \ru{с} \ru{ремонтом}.
\end{quote}

\begin{quote}
\emph{Svyókor pomóg nam s remóntom.}
\end{quote}

\begin{quote}
My father-in-law helped us with the renovation.
\end{quote}

Masculine; genitive \ru{свёкра} (the vowel drops). The husband's father.

\begin{cousinweb}
\emph{\ru{свёкор}} is inherited from Proto-Slavic \emph{svekr\ru{ъ}} through Old East Slavic \emph{svekr\ru{ъ}}, which EtymDB traces to PIE \emph{swé\'{k}uros} 'father-in-law'; \emph{\ru{свекровь}} goes back to the related PIE \emph{swe\'{k}rúh\textsubscript{2}} (EtymDB).
\end{cousinweb}

\begin{grammarlens}[title={What you've built}]
One more word for telling, explaining and describing at B1.
\end{grammarlens}

\subsection*{Your turn}
\begin{itemize}
\raggedright
  \item \emph{Say it:} \emph{svyókor}
  \item \emph{Say it:} \emph{svyókor}, once more
  \item \emph{Say it:} a sentence of your own with \emph{svyókor}
  \item \emph{Recall:} say the Russian for a husband's mother, then the Russian for a wife's mother, then say \emph{svyókor} again
\end{itemize}

\begin{tcolorbox}[breakable,colback=teal!4,colframe=teal!35!black,title={Before you move on}]
What does \ru{свёкор} mean? (\textquotedblleft{}A husband's father, a father-in-law\textquotedblright{}.) Say it once more.
\end{tcolorbox}
\section[\texorpdfstring{\ru{тесть} (test)}{тесть (test)}]{\ru{тесть} (test) — a wife's father, a father-in-law}

\label{lesson:RU-C491-test}



\begin{quote}
Before the new one: say the Russian for a wife's mother, then the Russian for a husband's father.
\end{quote}

\subsection*{You'll want to know: \ru{тесть}}
\textbf{\ru{тесть}} (\emph{test}) — \textquotedblleft{}a wife's father, a father-in-law\textquotedblright{}.

\begin{quote}
\ru{Мой} \ru{тесть} — \ru{бывший} \ru{лётчик}.
\end{quote}

\begin{quote}
\emph{Moy test — bývshiy lyótchik.}
\end{quote}

\begin{quote}
My father-in-law is a former pilot.
\end{quote}

Masculine. The wife's father.

\begin{cousinweb}
\emph{\ru{тесть}} is inherited from Proto-Slavic \emph{t\ru{ь}st\ru{ь}} (EtymDB).
\end{cousinweb}

\begin{grammarlens}[title={What you've built}]
One more word for telling, explaining and describing at B1.
\end{grammarlens}

\subsection*{Your turn}
\begin{itemize}
\raggedright
  \item \emph{Say it:} \emph{test}
  \item \emph{Say it:} \emph{test}, once more
  \item \emph{Say it:} a sentence of your own with \emph{test}
  \item \emph{Recall:} say the Russian for a wife's mother, then the Russian for a husband's father, then say \emph{test} again
\end{itemize}

\begin{tcolorbox}[breakable,colback=teal!4,colframe=teal!35!black,title={Before you move on}]
What does \ru{тесть} mean? (\textquotedblleft{}A wife's father, a father-in-law\textquotedblright{}.) Say it once more.
\end{tcolorbox}
\section[\texorpdfstring{\ru{зять} (zyat)}{зять (zyat)}]{\ru{зять} (zyat) — a son-in-law; a sister's husband}

\label{lesson:RU-C491-zyat}



\begin{quote}
Before the new one: say the Russian for a husband's father, then the Russian for a wife's father.
\end{quote}

\subsection*{You'll want to know: \ru{зять}}
\textbf{\ru{зять}} (\emph{zyat}) — \textquotedblleft{}a son-in-law; a sister's husband\textquotedblright{}.

\begin{quote}
\ru{На} \ru{свадьбе} \ru{тёща} \ru{танцевала} \ru{с} \ru{зятем}.
\end{quote}

\begin{quote}
\emph{Na svádbe tyóshcha tantsevála s zyátem.}
\end{quote}

\begin{quote}
At the wedding the mother-in-law danced with her son-in-law.
\end{quote}

Masculine; plural \ru{зятья}.

\begin{cousinweb}
\emph{\ru{зять}} is inherited from Proto-Slavic \emph{zęt\ru{ь}}, which EtymDB traces to PIE \emph{ǵenh\textsubscript{1}- 'to beget'.}
\end{cousinweb}

\begin{grammarlens}[title={What you've built}]
One more word for telling, explaining and describing at B1.
\end{grammarlens}

\subsection*{Your turn}
\begin{itemize}
\raggedright
  \item \emph{Say it:} \emph{zyat}
  \item \emph{Say it:} \emph{zyat}, once more
  \item \emph{Say it:} a sentence of your own with \emph{zyat}
  \item \emph{Recall:} say the Russian for a husband's father, then the Russian for a wife's father, then say \emph{zyat} again
\end{itemize}

\begin{tcolorbox}[breakable,colback=teal!4,colframe=teal!35!black,title={Before you move on}]
What does \ru{зять} mean? (\textquotedblleft{}A son-in-law; a sister's husband\textquotedblright{}.) Say it once more.
\end{tcolorbox}
